%% kapa: exam=1 round=33 year=2022 period=1 form=A subject=civil_law subject_ko=민법 seq=13 number=13
%% source: Q-Net 감정평가사 기출문제 (cst00309) artlSeq=5209997 file=raw/1008/5209997/2022년도 제33회 감정평가사 1차시험 1교시(민법,경제학,부동산학).pdf
%% conversion: kordoc 4.24.51 -> tools/gen1.py (auto); review=proofread-sonnet; flags=none
\begin{question}{13}
상대방 있는 의사표시의 효력발생에 관한 설명으로 옳은 것은? (다툼이 있으면 판례에 따름)
\begin{choices}
  \item 의사표시의 도달은 표의자의 상대방이 이를 현실적으로 수령하거나 그 통지의 내용을 알았을 것을 요한다.
  \item 제한능력자는 원칙적으로 의사표시의 수령무능력자이다.
  \item 보통우편의 방법으로 발송된 의사표시는 상당기간 내에 도달하였다고 추정된다.
  \item 표의자가 의사표시를 발송한 후 사망한 경우, 그 의사표시는 효력을 잃는다.
  \item 표의자가 과실로 상대방을 알지 못하는 경우에는 민사소송법 공시송달 규정에 의하여 의사표시의 효력을 발생시킬 수 있다.
\end{choices}
\end{question}
\answer{2}
